\documentclass{report}
\usepackage{amsmath,amssymb}
\setlength{\parindent}{0mm}
\setlength{\parskip}{1em}
\begin{document}
\begin{center}
\rule{6in}{1pt} \
{\large Andrey Prokopenko \\
{\bf Speeding up multi-physics simulations through reuse of multigrid components}}

Sandia National Laboratories \\ MailStop 1318 \\ PO BOX 5800 \\ Albuquerque \\ NM 87185-1318
\\
{\tt aprokop@sandia.gov}\\
Paul Lin\\
John Shadid\\
Jonathan Hu\end{center}

Large scale scientific applications, such as fluid dynamics or MHD, rely
on robust and efficient linear and nonlinear solvers for their
simulations. The complex nature of such simulations may require taking
small time steps which may result in a long time to solution, a
significant part of which is taken by linear solvers. The sheer amount of
work and transient nature of these problems allow for various approaches
to reduce the overall simulation time.
For example, nonlinear solvers for each transient step often produce a
sequence of closely related systems, and a promising strategy is to
leverage any such system similarities from within the linear solve phase.

In this talk, we will focus on simulations that use algebraic multigrid
(AMG) methods as part of their solution process. We will concentrate on
speeding up the AMG solver through reuse of
components of the multigrid hierarchy throughout the nonlinear solve, and
possibly across transient steps. The goal is to significantly reduce the
multigrid hierarchy setup time while not
significantly hurting the convergence. We will consider several AMG reuse
strategies that differ in the amount of reused data. We will demonstrate
the effectiveness of such strategies in parallel
applications.


\end{document}
